% =========================
% Cas pratique 10
% =========================

\section[Cas 10]{Cas pratique 10 : Contrefaçon
par fourniture de moyens}

%--------------------------
\begin{frame}{Question}

\begin{block}{}
L'article \textbf{26} AJBU (droit d'empêcher d'exploitation
indirecte de l'invention).
\end{block}

\begin{longtblr}{colspec = {X[1,l] X[3,l] X[3,l]}, rowsep = 0pt}

Critère & Contrefaçon & Pas contrefaçon \\

\hline
territoralité
& Papgeneric livre et offre de livrer sur
le territoire des Etats membres contractants
(implicite)
& \\

moyens essentiels
& du point de vue juridique oui, la revendication 1
comprend cette caractéristique
& du point de vue
technique non, l'invention fonctionne sans papier
avec des coins spéciaux dans un mode dégradé \\

personne habilité
& une personne autre que celle habilitée
à exploiter l'invention brevetée (implicite).
& \\

moyens aptes et destinés
& moyens aptes oui 
& moyens destinés difficile à vérifier. \\

\end{longtblr}

\end{frame}